\documentclass[UTF8,11pt]{ctexart}
\usepackage{amsmath,amssymb,geometry}
\geometry{margin=22mm}
\begin{document}
\section*{P49: 噪声感知量子位配置的加强线性整数模型}
\paragraph{来源与范围}依据原论文PDF第1, 2, 3, 4, 10, 11, 12, 13, 14, 15, 16页，第3节原式(1b)–(1k)、(12)–(17)；串扰按原式(2)同型McCormick线性化。

\subsection*{集合、参数与变量}
$Q$为逻辑量子位集合；$H=(V,A_H)$为连通的硬件有向图，$|V|\ge|Q|$，且弧集合对称。
$N\in\mathbb Z_{++}$为电路层数，$L^t\subseteq Q\times Q$为第$t$层两两不共享逻辑量子位的双比特门。
$d_H(i,j)$为硬件图最短路距离，$\lambda\in[0,1]$为给定权衡系数。
$E_i\ge0$为单比特门/读出噪声费用，$E_{i,j}\ge0$为有向双比特门噪声费用，均需对使用的硬件元素明确给出。
令$N_H(i)=\{j:(i,j)\in A_H\}$，$N((i,j))=(N_H(i)\setminus\{j\})\cup(N_H(j)\setminus\{i\})$。
给定实测错误率比$r_{(i,j)}^k\ge0$，串扰系数为$E_{(i,j),k}=\max\{0,(r_{(i,j)}^k-1)E_k\}$。
变量$w_{q,i}^t$表示层$t$逻辑位$q$占据硬件位$i$；$x_{q,i,j}^t$表示$q$在相邻层间从$i$转移至$j$；
$z_{(p,q),(i,j)}^t$表示逻辑门$(p,q)$在硬件弧$(i,j)$执行；$u_{q,(r,s),(i,j),k}^t$为$w_{q,k}^tz_{(r,s),(i,j)}^t$的线性化变量。
$u$的完整索引域是$1\le t\le N,(r,s)\in L^t,(i,j)\in A_H,k\in N((i,j)),q\in Q$。
初始布局由模型自由优化，所有层使用相同硬件子图。以下$1a$等为本整理编号，来源是原文(1)、(12)–(17)。
\subsection*{噪声感知TAP的完整线性整数形式}
\begin{equation}
\min \;
(1-\lambda)c_{\mathrm{swap}} + \lambda c_{\mathrm{noise}}
\tag{1a}
\end{equation}

\begin{equation}
c_{\mathrm{swap}} =
\sum_{t=1}^{N-1} \sum_{q\in Q} \sum_{i,j\in V\times V}
d_H(i,j)x_{q,i,j}^t
\tag{1b}
\end{equation}

\begin{equation}
c_{\mathrm{noise}} =
\sum_{i\in V} E_i \sum_{t=1}^N \sum_{q\in Q} w_{q,i}^t
+
\sum_{(i,j)\in A_H} E_{i,j} \sum_{t=1}^N \sum_{(p,q)\in L^t} z_{(p,q),(i,j)}^t
+
\sum_{(i,j)\in A_H} \sum_{k\in N((i,j))} E_{(i,j),k}
\sum_{t=1}^N \sum_{(r,s)\in L^t} \sum_{q\in Q}
u_{q,(r,s),(i,j),k}^t
\tag{1c}
\end{equation}

\begin{equation}
w_{q,i}^t = \sum_{j\in V} x_{q,i,j}^t,
\qquad
\forall 1 \le t \le N-1,\ \forall i \in V,\ \forall q \in Q
\tag{1d}
\end{equation}

\begin{equation}
w_{q,i}^{t+1} = \sum_{j\in V} x_{q,j,i}^t,
\qquad
\forall 1 \le t \le N-1,\ \forall i \in V,\ \forall q \in Q
\tag{1e}
\end{equation}

\begin{equation}
\sum_{i\in V} w_{q,i}^t = 1,
\qquad
\forall 1 \le t \le N,\ \forall q \in Q
\tag{1f}
\end{equation}

\begin{equation}
\sum_{q\in Q} w_{q,i}^t \le 1,
\qquad
\forall 1 \le t \le N,\ \forall i \in V
\tag{1g}
\end{equation}

\begin{equation}
\sum_{(i,j)\in A_H} z_{(p,q),(i,j)}^t = 1,
\qquad
\forall 1 \le t \le N,\ \forall (p,q) \in L^t
\tag{1h}
\end{equation}



\begin{equation}
\sum_{q\in Q} w_{q,i}^t = \sum_{q\in Q} w_{q,i}^1,
\qquad
\forall 2 \le t \le N,\ \forall i \in V
\tag{1j}
\end{equation}

\begin{equation}
\sum_{j:(i,j)\in A_H} z_{(p,q),(i,j)}^t = w_{p,i}^t,
\qquad
\forall 1 \le t \le N,\ \forall (p,q) \in L^t,\ \forall i \in V
\tag{2a}
\end{equation}

\begin{equation}
\sum_{i:(i,j)\in A_H} z_{(p,q),(i,j)}^t = w_{q,j}^t,
\qquad
\forall 1 \le t \le N,\ \forall (p,q) \in L^t,\ \forall j \in V
\tag{2b}
\end{equation}

\begin{equation}
u_{q,(r,s),(i,j),k}^t \le w_{q,k}^t,
\qquad
\forall t,(r,s),(i,j),k,q
\tag{2c}
\end{equation}

\begin{equation}
u_{q,(r,s),(i,j),k}^t \le z_{(r,s),(i,j)}^t,
\qquad
\forall t,(r,s),(i,j),k,q
\tag{2d}
\end{equation}

\begin{equation}
u_{q,(r,s),(i,j),k}^t \ge w_{q,k}^t + z_{(r,s),(i,j)}^t - 1,
\qquad
\forall t,(r,s),(i,j),k,q
\tag{2e}
\end{equation}

\begin{equation}
w_{q,i}^t \in \{0,1\},
\qquad
\forall 1 \le t \le N,\ \forall q \in Q,\ \forall i \in V
\tag{3a}
\end{equation}

\begin{equation}
x_{q,i,j}^t \in \{0,1\},
\qquad
\forall 1 \le t < N,\ \forall q \in Q,\ \forall (i,j) \in V \times V
\tag{3b}
\end{equation}

\begin{equation}
z_{(p,q),(i,j)}^t \in \{0,1\},
\qquad
\forall 1 \le t \le N,\ \forall (p,q) \in L^t,\ \forall (i,j) \in A_H
\tag{3c}
\end{equation}

\begin{equation}
u_{q,(r,s),(i,j),k}^t \in \{0,1\},
\qquad
\forall t,(r,s),(i,j),k,q
\tag{3d}
\end{equation}


\end{document}
